\begin{abstract}
把遥测数据留在内部的组织，只能用小型本地 LLM 来分诊告警，而这些模型单独做不到。不过，许多告警属于反复出现的类型：分析员已经知道它们的候选原因、区分这些原因的检查，以及以往案例的结论。我们提出 \hesp{}，通过把分诊拆成两项工作，让小模型能用于这类告警：模型负责\emph{读}，把每个探针返回的原始日志文本读成分析员定义好的某个结果；控制器负责\emph{决定}下一步运行哪个探针、何时停止、接受哪个结论，依据的是候选原因上的贝叶斯账本，其似然从已结案的案例中统计得到。在三个模拟分诊场景中，用五个 7B 到 72B 的开源权重模型测试：小模型自行决定时，即使每条观测都已给好标签，最多也只解决三分之一的案例；必须自己读原始日志时，48 个案例一个也没有解决。观测已结构化时，控制器单独就解决了其中两个场景的全部案例；但日志变更改变了日志格式后，它的规则解析器读不出 364 条观测中的任何一条，每个案例都被升级。两者结合，由 Qwen2.5-7B 读取、控制器决定时，48 个漂移格式回合全部解决，Llama-3.1-8B 解决 34 个。读取也是攻击通道：攻击者把一致的良性故事写进每一条日志，会让作为读取器的 Qwen2.5-7B 把全部 36 个攻击结为良性。永远不让模型解析出的证据支持良性结论，挡住了我们测试的每一种攻击，代价是在没有可信解析器能读日志时，诚实的良性案例会被升级。我们的场景是模拟的，其中的原因与结果是事先写好的，且每个原因都会留下一条独特的观测；新的告警类型和从真实工单统计的表不在其范围之内。
\end{abstract}
